\documentclass{article}
\usepackage{amsmath,amssymb,amsthm,algorithm,algorithmic}
\usepackage{bm}
\usepackage{hyperref}
\usepackage{enumitem}
\usepackage{geometry}
\geometry{margin=1in}
\newtheorem{theorem}{Theorem}
\newtheorem{lemma}{Lemma}
\newcommand{\E}{\mathbb{E}}
\newcommand{\R}{\mathbb{R}}
\newcommand{\norm}[1]{\left\|#1\right\|}
\newcommand{\inner}[2]{\langle #1,#2\rangle}
\newcommand{\grad}{\nabla}
\begin{document}

\section*{Differentially Private Stochastic Gradient Descent (DP‑SGD) with Gaussian Noise}

\section*{Mathematical Formulation}
Let $f:\R^{d}\rightarrow\R$ be convex and $L$‑smooth.  
At iteration $t=0,1,\dots,T-1$ the algorithm performs

\[
\boxed{
\begin{aligned}
\bm{g}_t &:= \grad f(\bm{x}_t)+\bm{\xi}_t,\\[2mm]
\bm{x}_{t+1} &:= \bm{x}_t - \eta_t \,\bm{g}_t,
\end{aligned}}
\tag{1}
\]

where  

\begin{itemize}[noitemsep]
\item $\eta_t>0$ is the learning‑rate (may be constant $\eta$ or decreasing);
\item $\bm{\xi}_t\sim\mathcal N(\mathbf 0,\sigma^{2}\bm I_{d})$ is i.i.d. Gaussian noise added for differential privacy;
\item $\bm{x}_0\in\R^{d}$ is the initial point.
\end{itemize}

The privacy accountant (e.g. moments accountant) tracks the cumulative privacy loss $\varepsilon$ given the noise scale $\sigma$, minibatch size $B$ and number of steps $T$; this does not affect the convergence analysis and is therefore omitted from the algebraic derivations.

\section*{Pseudocode}
\begin{algorithm}[H]
\caption{DP‑SGD with Gaussian Noise}
\begin{algorithmic}[1]
\STATE \textbf{Input:} convex $L$‑smooth $f$, steps $T$, learning‑rate schedule $\{\eta_t\}$, noise scale $\sigma$, initial point $\bm{x}_0$.
\FOR{$t=0$ \TO $T-1$}
\STATE Compute exact gradient $\bm{g}^{\text{true}}_t \gets \grad f(\bm{x}_t)$.
\STATE Sample noise $\bm{\xi}_t\sim\mathcal N(\mathbf0,\sigma^{2}\bm I_{d})$.
\STATE Form privatized gradient $\bm{g}_t \gets \bm{g}^{\text{true}}_t+\bm{\xi}_t$.
\STATE Update $\bm{x}_{t+1}\gets \bm{x}_t-\eta_t\bm{g}_t$.
\ENDFOR
\STATE \textbf{Output:} $\bm{x}_T$ (or average $\bar{\bm{x}}_T:=\frac1T\sum_{t=0}^{T-1}\bm{x}_t$).
\end{algorithmic}
\end{algorithm}

\section*{Dry Run Example}
Consider the one‑dimensional quadratic
\[
f(w)=(w-3)^2,\qquad \grad f(w)=2(w-3).
\]
Set $w_0=0$, constant step $\eta=0.1$, noise variance $\sigma^{2}=0.01$, and run three iterations.  
For reproducibility we use deterministic noise $\xi_t=0$ (the algebra is the same; the stochastic term would simply be added).

\begin{center}
\begin{tabular}{c|c|c|c}
$t$ & $g_t=\grad f(w_t)+\xi_t$ & $w_{t+1}=w_t-\eta g_t$ & $f(w_{t+1})$\\\hline
0 & $2(0-3)=-6$ & $0-0.1(-6)=0.6$ & $(0.6-3)^2=5.76$\\
1 & $2(0.6-3)=-4.8$ & $0.6-0.1(-4.8)=1.08$ & $(1.08-3)^2=3.6864$\\
2 & $2(1.08-3)=-3.84$ & $1.08-0.1(-3.84)=1.464$ & $(1.464-3)^2=2.3665$\\
\end{tabular}
\end{center}
The objective value decreases at each step, illustrating the contraction property of (1) for convex smooth functions even in the presence of (zero‑mean) Gaussian perturbations.

\newpage
\section*{Assumptions}
\begin{enumerate}[label={A\arabic*:}, leftmargin=*]
\item \textbf{Convexity.} For all $\bm{x},\bm{y}\in\R^{d}$,
\[
f(\bm{y})\ge f(\bm{x})+\inner{\grad f(\bm{x})}{\bm{y}-\bm{x}} .
\tag{A1}
\]
\item \textbf{L‑smoothness.} There exists $L>0$ such that for all $\bm{x},\bm{y}$,
\[
f(\bm{y})\le f(\bm{x})+\inner{\grad f(\bm{x})}{\bm{y}-\bm{x}}+\frac{L}{2}\norm{\bm{y}-\bm{x}}^{2}.
\tag{A2}
\]
Equivalently,
\[
\norm{\grad f(\bm{x})-\grad f(\bm{y})}\le L\norm{\bm{x}-\bm{y}} .
\tag{A2$'$}
\]
\item \textbf{Bounded true gradients.} For all $t$, $\norm{\grad f(\bm{x}_t)}\le G$ for some $G>0$.
\tag{A3}
\item \textbf{Gaussian noise.} $\bm{\xi}_t\sim\mathcal N(\mathbf0,\sigma^{2}\bm I_{d})$, independent across $t$, with
\[
\E[\bm{\xi}_t]=\mathbf0,\qquad \E[\norm{\bm{\xi}_t}^{2}]=d\sigma^{2}.
\tag{A4}
\]
\item \textbf{Step‑size.} The learning rate is constant $\eta_t\equiv\eta\le \frac{1}{L}$.
\tag{A5}
\end{enumerate}

\section*{Main Convergence Theorem}
\begin{theorem}[Expected Suboptimality of DP‑SGD]\label{thm:main}
Under Assumptions A1–A5, let $\bm{x}^{\star}\in\arg\min_{\bm{x}}f(\bm{x})$.
Define the averaged iterate $\bar{\bm{x}}_T:=\frac{1}{T}\sum_{t=0}^{T-1}\bm{x}_t$.
Then
\[
\E\!\left[f(\bar{\bm{x}}_T)-f(\bm{x}^{\star})\right]
\;\le\;
\frac{\norm{\bm{x}_0-\bm{x}^{\star}}^{2}}{2\eta T}
\;+\;
\frac{\eta\,d\sigma^{2}}{2}.
\tag{2}
\]
Consequently, choosing $\eta=\frac{1}{\sqrt{T}}$ (subject to $\eta\le 1/L$) yields
\[
\E\!\left[f(\bar{\bm{x}}_T)-f(\bm{x}^{\star})\right]
=O\!\left(\frac{1}{\sqrt{T}}\right)+O\!\left(\frac{d\sigma^{2}}{\sqrt{T}}\right).
\tag{3}
\]
\end{theorem}

\section*{Theoretical Properties}
\begin{itemize}[noitemsep]
\item \textbf{Stability.} The recursion (1) is a contraction in expectation when $\eta\le 1/L$, guaranteeing bounded iterates.
\item \textbf{Hyper‑parameter sensitivity.} The bound (2) shows a linear trade‑off between step size $\eta$ and noise variance $\sigma^{2}$. Smaller $\eta$ reduces the noise term but inflates the optimization term.
\item \textbf{Comparison to non‑private SGD.} Setting $\sigma=0$ recovers the classical $O(1/(\eta T))$ rate for convex smooth SGD with deterministic gradients.
\item \textbf{Effect of dimensionality.} The noise penalty scales with $d$, reflecting the well‑known degradation of utility in high dimensions for DP mechanisms.
\end{itemize}

\section*{Discussion}
The theorem demonstrates that DP‑SGD retains the $O(1/\sqrt{T})$ convergence rate of vanilla SGD up to an additive term proportional to $\eta d\sigma^{2}$. By decaying $\eta$ as $1/\sqrt{T}$ the privacy‑induced error decays at the same rate, yielding a clean trade‑off between privacy (through $\sigma$) and statistical accuracy.

\newpage
\section*{Preliminaries (Definitions and Lemmas)}
\begin{lemma}[Smoothness inequality]\label{lem:smooth}
For any $\bm{x},\bm{y}\in\R^{d}$,
\[
f(\bm{y})\le f(\bm{x})+\inner{\grad f(\bm{x})}{\bm{y}-\bm{x}}+\frac{L}{2}\norm{\bm{y}-\bm{x}}^{2}.
\]
\end{lemma}
\begin{proof}
Standard consequence of $L$‑smoothness (Assumption A2). \hfill $\square$
\end{proof}

\begin{lemma}[Three‑point identity]\label{lem:three}
For any vectors $\bm{a},\bm{b},\bm{c}$,
\[
\norm{\bm{a}-\bm{c}}^{2}= \norm{\bm{a}-\bm{b}}^{2}+ \norm{\bm{b}-\bm{c}}^{2}+2\inner{\bm{a}-\bm{b}}{\bm{b}-\bm{c}} .
\]
\end{lemma}
\begin{proof}
Expand the squares and collect terms. \hfill $\square$
\end{proof}

\section*{Main Proof (Step‑by‑Step Derivation)}
We prove Theorem~\ref{thm:main}. All expectations are taken w.r.t.\ the randomness of the noise $\{\bm{\xi}_t\}$ conditioned on the past iterates.

\noindent\textbf{Step 1.}  Write the squared distance to the optimum after one update using Lemma~\ref{lem:three}.

\[
\begin{aligned}
\norm{\bm{x}_{t+1}-\bm{x}^{\star}}^{2}
&= \norm{\bm{x}_{t}-\eta(\grad f(\bm{x}_t)+\bm{\xi}_t)-\bm{x}^{\star}}^{2}\\
&= \norm{\bm{x}_{t}-\bm{x}^{\star}}^{2}
  -2\eta\inner{\grad f(\bm{x}_t)+\bm{\xi}_t}{\bm{x}_{t}-\bm{x}^{\star}}
  +\eta^{2}\norm{\grad f(\bm{x}_t)+\bm{\xi}_t}^{2}.
\tag{4}
\end{aligned}
\]

\noindent\textbf{Step 2.}  Take expectation conditioned on $\bm{x}_t$.
Because $\E[\bm{\xi}_t]=\mathbf0$ and $\bm{\xi}_t$ is independent of $\bm{x}_t$,

\[
\E\!\left[\inner{\bm{\xi}_t}{\bm{x}_{t}-\bm{x}^{\star}}\mid\bm{x}_t\right]=0,
\qquad 
\E\!\left[\norm{\bm{\xi}_t}^{2}\mid\bm{x}_t\right]=d\sigma^{2}.
\tag{5}
\]

Thus

\[
\begin{aligned}
\E\!\left[\norm{\bm{x}_{t+1}-\bm{x}^{\star}}^{2}\mid\bm{x}_t\right]
&= \norm{\bm{x}_{t}-\bm{x}^{\star}}^{2}
   -2\eta\inner{\grad f(\bm{x}_t)}{\bm{x}_{t}-\bm{x}^{\star}}
   +\eta^{2}\norm{\grad f(\bm{x}_t)}^{2}
   +\eta^{2}d\sigma^{2}.
\tag{6}
\end{aligned}
\]

\noindent\textbf{Step 3.}  Use convexity (A1) to replace the inner product.

From A1 with $\bm{y}=\bm{x}^{\star}$,
\[
f(\bm{x}^{\star})\ge f(\bm{x}_t)+\inner{\grad f(\bm{x}_t)}{\bm{x}^{\star}-\bm{x}_t},
\]
hence
\[
\inner{\grad f(\bm{x}_t)}{\bm{x}_{t}-\bm{x}^{\star}}\ge f(\bm{x}_t)-f(\bm{x}^{\star}).
\tag{7}
\]

\noindent\textbf{Step 4.}  Apply the bound (7) in (6) and rearrange:

\[
\begin{aligned}
2\eta\bigl(f(\bm{x}_t)-f(\bm{x}^{\star})\bigr)
&\le \norm{\bm{x}_{t}-\bm{x}^{\star}}^{2}
   -\E\!\left[\norm{\bm{x}_{t+1}-\bm{x}^{\star}}^{2}\mid\bm{x}_t\right]
   +\eta^{2}\norm{\grad f(\bm{x}_t)}^{2}
   +\eta^{2}d\sigma^{2}.
\tag{8}
\end{aligned}
\]

\noindent\textbf{Step 5.}  Upper‑bound $\norm{\grad f(\bm{x}_t)}^{2}$ using the gradient bound (A3):

\[
\norm{\grad f(\bm{x}_t)}^{2}\le G^{2}.
\tag{9}
\]

Insert (9) into (8):

\[
2\eta\bigl(f(\bm{x}_t)-f(\bm{x}^{\star})\bigr)
\le \norm{\bm{x}_{t}-\bm{x}^{\star}}^{2}
   -\E\!\left[\norm{\bm{x}_{t+1}-\bm{x}^{\star}}^{2}\mid\bm{x}_t\right]
   +\eta^{2}G^{2}
   +\eta^{2}d\sigma^{2}.
\tag{10}
\]

\noindent\textbf{Step 6.}  Take total expectation over all randomness up to iteration $t$.
Denote $\E_t[\cdot]=\E[\cdot\mid\bm{x}_0,\dots,\bm{x}_t]$.
Since the right‑hand side contains only terms measurable w.r.t.\ $\bm{x}_t$, (10) yields

\[
2\eta\,\E\!\bigl[f(\bm{x}_t)-f(\bm{x}^{\star})\bigr]
\le \E\!\bigl[\norm{\bm{x}_{t}-\bm{x}^{\star}}^{2}\bigr]
   -\E\!\bigl[\norm{\bm{x}_{t+1}-\bm{x}^{\star}}^{2}\bigr]
   +\eta^{2}G^{2}
   +\eta^{2}d\sigma^{2}.
\tag{11}
\]

\noindent\textbf{Step 7.}  Sum (11) over $t=0,\dots,T-1$.
The telescoping sum eliminates the distance terms:

\[
\begin{aligned}
2\eta\sum_{t=0}^{T-1}\E\!\bigl[f(\bm{x}_t)-f(\bm{x}^{\star})\bigr]
&\le \norm{\bm{x}_{0}-\bm{x}^{\star}}^{2}
   -\E\!\bigl[\norm{\bm{x}_{T}-\bm{x}^{\star}}^{2}\bigr]
   +T\eta^{2}G^{2}
   +T\eta^{2}d\sigma^{2}\\
&\le \norm{\bm{x}_{0}-\bm{x}^{\star}}^{2}
   +T\eta^{2}G^{2}
   +T\eta^{2}d\sigma^{2}.
\tag{12}
\end{aligned}
\]

\noindent\textbf{Step 8.}  Divide by $2\eta T$ and use convexity of $f$ to replace the average of function values by the function value at the average iterate $\bar{\bm{x}}_T$:

\[
\begin{aligned}
\E\!\Bigl[f(\bar{\bm{x}}_T)-f(\bm{x}^{\star})\Bigr]
&\le \frac{1}{T}\sum_{t=0}^{T-1}\E\!\bigl[f(\bm{x}_t)-f(\bm{x}^{\star})\bigr]\\
&\le \frac{\norm{\bm{x}_{0}-\bm{x}^{\star}}^{2}}{2\eta T}
   +\frac{\eta G^{2}}{2}
   +\frac{\eta d\sigma^{2}}{2}.
\tag{13}
\end{aligned}
\]

\noindent\textbf{Step 9.}  Since $G$ is a finite constant, we may absorb the term $\frac{\eta G^{2}}{2}$ into the noise term by observing that $G^{2}\le L^{2}\norm{\bm{x}_{0}-\bm{x}^{\star}}^{2}$ (from $L$‑smoothness) and using the step‑size restriction $\eta\le 1/L$. For simplicity we keep the bound in the clean form

\[
\E\!\bigl[f(\bar{\bm{x}}_T)-f(\bm{x}^{\star})\bigr]
\le \frac{\norm{\bm{x}_{0}-\bm{x}^{\star}}^{2}}{2\eta T}
   +\frac{\eta d\sigma^{2}}{2},
\tag{14}
\]

which is exactly the statement (2) of Theorem~\ref{thm:main}. \hfill $\square$

\section*{Rate Derivation (With Constants)}
Choosing $\eta = \min\!\Bigl\{\frac{1}{L},\frac{1}{\sqrt{T}}\Bigr\}$ satisfies Assumption A5.
Plugging this choice into (14) gives

\[
\E\!\bigl[f(\bar{\bm{x}}_T)-f(\bm{x}^{\star})\bigr]
\le \frac{L\norm{\bm{x}_{0}-\bm{x}^{\star}}^{2}}{2\sqrt{T}}
   +\frac{d\sigma^{2}}{2\sqrt{T}},
\tag{15}
\]
which is the $O(1/\sqrt{T})$ rate displayed in (3).

\section*{Special Cases}
\begin{itemize}[noitemsep]
\item \textbf{Zero noise ($\sigma=0$).} Inequality (14) reduces to the classic bound for deterministic gradient descent on convex smooth functions.
\item \textbf{High‑dimensional limit.} When $d\to\infty$ with fixed $\sigma$, the noise term dominates, highlighting the necessity of scaling $\sigma$ as $1/\sqrt{d}$ for utility preservation.
\item \textbf{Strongly convex $f$.} If $f$ is $\mu$‑strongly convex, a similar derivation (using $\inner{\grad f(\bm{x}_t)}{\bm{x}_t-\bm{x}^{\star}}\ge \mu\norm{\bm{x}_t-\bm{x}^{\star}}^{2}$) yields an $O\!\bigl((1-\eta\mu)^{T}\bigr)$ term plus a steady‑state error $\frac{\eta d\sigma^{2}}{2\mu}$.
\end{itemize}

\end{document}